% === Revised Section 8: Caveats ===
%
% NEW PREAMBLE ADDITIONS NEEDED:
% ------------------------------------------------------------------
% \usepackage[most]{tcolorbox}   % already in main preamble
% \tcbuselibrary{breakable,skins} % already in main preamble
%
% \newtcolorbox{activity}[1][]{
%   breakable,
%   enhanced,
%   colback=blue!5,
%   colframe=blue!55!black,
%   coltitle=white,
%   fonttitle=\bfseries,
%   title={Activity},
%   boxrule=0.6pt,
%   left=6pt,right=6pt,top=4pt,bottom=4pt,
%   #1
% }
%
% \newtcolorbox{procurement}[1][]{
%   breakable,
%   enhanced,
%   colback=green!5,
%   colframe=green!45!black,
%   coltitle=white,
%   fonttitle=\bfseries,
%   title={Procurement corollary},
%   boxrule=0.5pt,
%   left=6pt,right=6pt,top=3pt,bottom=3pt,
%   fontupper=\small,
%   #1
% }
% ------------------------------------------------------------------

\section{Caveats and what the equivalence does not say}
\label{sec:caveats}

\begin{intuition}
The equivalence is real but narrow. It tells you that good prediction
entails good \emph{modelling}; it does \emph{not} tell you that good
prediction entails understanding, agency, or safety. It promises
statistical competence. It does not promise causal reasoning,
counterfactual simulation, or moral judgement. Those may emerge ---
the theory leaves the door open --- but they are not guaranteed by
the objective.

\medskip

\noindent\textbf{The food-truck and the Michelin kitchen.} Keep this
image in your pocket for the rest of the section. Uncomputability is
not cosmetic: $K(x)$ and AIXI are unreachable
(\cref{thm:uncomputable}), and every real compressor, every real
agent, is a finite approximation. Saying an LLM is
``approximating AIXI'' is like a food-truck claiming to approximate
a Michelin three-star kitchen by reading the menu. Technically on
the same axis. In practice, not the relevant comparison. The rest
of this section is a tour of five specific places where the axis
and the kitchen come apart.

\medskip

\noindent\textbf{The running cast.} To keep the five abstract
caveats walkable, we will reuse three characters:

\begin{itemize}[leftmargin=1.2em,itemsep=2pt]
  \item \emph{Anna the junior analyst,} who has memorised every deal
        memo in the vault --- $4{,}000$ deals going back a decade.
  \item \emph{Ben the junior analyst,} who has read only ten of those
        memos but understands why each deal priced where it did.
  \item \emph{The phrasebook chatbot,} which has ``learned'' French by
        memorising 100{,}000 tourist phrases and can now order
        croissants flawlessly in thirteen dialects.
\end{itemize}

You will meet each of them again.

\medskip

\noindent\textbf{Remark 1 (the only non-tautological falsifier).}
A reasonable operator will ask: if AIXI is uncomputable and every
deployed system is ``a finite approximation,'' what test \emph{other
than cross-entropy itself} decides whether system A approximates the
ideal better than system B? The honest answer is: held-out
cross-entropy on distributions the training set did not see, plus
downstream task transfer. There is no deeper falsifier in the
current theory. If that makes the AIXI framing feel tautological,
the operator is reading the field correctly, not mis-reading the
theorem. The operational response is not philosophical (``but
really, AIXI!''). It is empirical: hold out a genuinely unseen
distribution, measure transfer, and stop there.

\medskip

\noindent\textbf{Compression $\ne$ understanding.}
Ask Anna, with her $4{,}000$ memorised memos, to price a deal that
is \emph{not} in the vault. What you learn from her answer is
whether she compressed the memos or the market. Ask the phrasebook
chatbot a sentence that is off-phrasebook --- ``my croissant has
opinions about the Maastricht Treaty'' --- and you find out whether
it has learned French or merely tiled the phrase space. A lossless
compressor has captured statistics, which is \emph{some} kind of
understanding, but not necessarily the kind a board means when it
asks ``does the model understand our business?''

\medskip

\noindent\textbf{Embodiment.}
A text-only system has only ever been surprised by text. It has no
ground-truth feedback from physical reality. Solomonoff's framework
accepts any computable environment, but LLMs have the text channel
and lack the body-in-a-world channel. Agentic, tool-using,
multimodal systems are early attempts to close the gap. The
equivalence does not tell you the gap has been closed.

\medskip

\noindent\textbf{Benchmarks are not the floor.}
A model with low cross-entropy on a benchmark may just have
memorised the benchmark. Data contamination is rampant. The
equivalence only holds when the test distribution is genuinely held
out --- which, for a model trained on ``the internet'', is an
increasingly expensive proposition.

\medskip

\noindent\textbf{Why a decision-maker should care.} Use the
equivalence as a \emph{lower bound, not an upper bound.} When a
vendor says ``lowest perplexity,'' accept that as real evidence of
genuine modelling skill. Then ask the second question: skill at
what, tested on what, does it transfer to your domain? The
equivalence tells you the knob exists. It does not tell you that
turning the knob solves your problem. Confusing ``next-token
prediction'' with ``just autocomplete'' misses a deep result;
confusing ``next-token prediction'' with ``general intelligence''
overclaims a narrow one. Both errors are expensive. Staying in the
middle --- taking the equivalence seriously but not theologically
--- is the managerial win.
\end{intuition}

\subsection*{Formal caveats (C1)--(C5)}
\label{sec:formal-caveats}

Several limitations of the chain~\eqref{eq:chain} deserve formal
acknowledgement. Each is followed by a one-line \emph{procurement
corollary}: the thing to actually say when a vendor pitches you.

\paragraph{(C1) Uncomputability.}
By \cref{thm:uncomputable}, $K(x)$ is not computable; by a related
argument so is $M(x)$. The entire right-hand side of~\eqref{eq:chain}
therefore names a mathematical \emph{ideal}, not an algorithm. Every
real compressor produces an upper bound $\tilde K(x) \ge K(x)$, and
equality is in principle unattainable. ``Approximates the universal
prior'' must be taken as an asymptotic statement about cross-entropy,
not as a literal claim that the model is computing $M$.

\begin{procurement}
If a vendor says their model ``approximates AIXI'' or ``computes the
Solomonoff prior,'' ask on what computable sub-class, with what
error bound, and over what distribution. The honest answer is
``we lower cross-entropy on held-out data.'' That is fine; it is
just not the same claim.
\end{procurement}

\paragraph{(C2) The additive constant.}
\Cref{thm:invariance} only guarantees machine-independence up to an
additive constant $c_{U_1, U_2}$. For short strings this constant
can dominate $K(x)$ itself; Kolmogorov complexity is only meaningful
asymptotically. Analogously, the $K(\mu)$ term
in~\eqref{eq:sol-bound} is a finite one-time information cost: it
does not shrink with data.

\begin{procurement}
On small, bespoke datasets (your legal corpus, your IoT telemetry)
the constants swamp the asymptotics. ``Big model + lots of data''
is the regime in which the theory speaks; ``small model + your
niche'' is the regime in which you hire evaluators, not theorists.
\end{procurement}

\paragraph{(C3) Practical compressors vs.\ $K$.}
A compressor $C$ is \emph{universal} in Shannon's sense if its code
length $L_C(x_{1:T})/T$ converges to the entropy rate of any
stationary ergodic source. It is \emph{not} in general universal in
Solomonoff's sense: it need not dominate all computable measures.
The gap between Lempel--Ziv (Shannon-universal) and Solomonoff
(universal over all computable measures) is a full logical level of
universality, and neural compressors live somewhere in between.

\begin{procurement}
``Compresses well'' comes in strengths. Ask which strength. A
model that crushes Wikipedia perplexity may still be a
Shannon-universal compressor on stationary text, not a Solomonoff-
universal one on your domain's non-stationarities.
\end{procurement}

\paragraph{(C4) Necessary vs.\ sufficient.}
\Cref{thm:solomonoff,thm:aixi-opt} establish that optimal prediction
and agency \emph{entail} optimal compression. They do \emph{not}
establish that every improvement in compression corresponds to an
improvement in every intuitive sense of ``understanding.'' In
particular, the framework is silent on phenomenal consciousness,
grounding, and embodiment, and it takes the reward signal and the
action--percept loop as primitives rather than explaining their
origin.

\begin{procurement}
Compression is \emph{necessary} for intelligence in the theory's
sense; no theorem says it is \emph{sufficient} for intelligence in
your sense. If the vendor's pitch relies on sufficiency, the
theory is not their friend.
\end{procurement}

\paragraph{(C5) Compression $\ne$ understanding (as internal structure).}
A compressor is characterised by its length function, not by its
internal representations. Two compressors with identical code
lengths may use entirely different internal mechanisms, only one of
which we would intuitively call ``understanding.''
Equation~\eqref{eq:chain} reduces ``intelligence'' to an extensional
(input--output) property; any argument that understanding is
intensional (about how the function is computed) is not touched by
the chain. Anna and Ben, from the running cast, can have identical
cross-entropy on the ten deals Ben has actually seen. They will
diverge on the eleventh.

\begin{procurement}
If you care \emph{how} the model arrives at its answers (audit,
explainability, regulation), cross-entropy is silent on your
question. You need mechanistic evaluation, not just a leaderboard
score.
\end{procurement}

\paragraph{Summary.}
The equivalence~\eqref{eq:chain} is a genuine theorem modulo
\crefrange{def:entropy}{def:aixi} and the computability-based
caveats (C1)--(C3), which are mathematical. Caveats (C4)--(C5) are
of a different species: they are not theorems against the chain,
they are \emph{scope restrictions} on what the chain claims.
What~\eqref{eq:chain} licences is a precise sense in which lowering
cross-entropy on a sufficiently rich corpus moves a predictor along
a single axis shared with compression and with Hutter-style
universal agency. What it does \emph{not} licence is the stronger
claim that this axis exhausts what we mean by understanding, safety,
or deployment-readiness.

% ----------------------------------------------------------------------
\subsection*{Activities}
\label{sec:caveats-activities}

\begin{activity}[title={Activity 8.1 --- Procurement checklist}]
You are briefing your CIO tomorrow on a vendor proposal. Derive six
questions --- one per caveat (C1)--(C5) plus one meta-question ---
to put to the vendor in writing. A suggested starter set:

\begin{enumerate}[leftmargin=1.4em,itemsep=2pt]
  \item (C1) ``On what benchmark, with what held-out distribution, is
        your perplexity number reported?''
  \item (C2) ``How was the model adapted to our domain? What is the
        performance delta between general and domain-specific
        held-out data?''
  \item (C3) ``Is your evaluation stationary or does it cover
        distribution shift? Give an example.''
  \item (C4) ``What evidence, beyond cross-entropy, supports the
        claim that the model \emph{understands} our data?''
  \item (C5) ``If we need to audit \emph{how} an answer was produced,
        what tooling do you provide?''
  \item (Meta) ``What would falsify your claims? Name the
        experiment.''
\end{enumerate}

\textbf{Deliverable.} A one-page checklist you could email to
procurement unchanged. Helps Aisha (a takeaway list) and Leo (the
formal caveats now have teeth).
\end{activity}

\begin{activity}[title={Activity 8.2 --- Benchmark-contamination detective}]
You are shown five short answers a model has given on a well-known
reasoning benchmark. Your job: flag which answers look
\emph{memorised} (verbatim phrasing, domain-specific jargon the
prompt did not contain, suspiciously specific numbers) versus
\emph{reasoned} (shows working, makes small errors, paraphrases).

For each flagged answer, write one sentence: what would you do next
to confirm contamination? (Suggested moves: perturb the prompt
wording; shift numerical values; ask a follow-up that is only
answerable from the derivation, not the final answer.)

\textbf{Deliverable.} A 5-row table: answer \#, flag
(memorised/reasoned/unsure), test-to-confirm. Helps Yuki (concrete
stimulus) and Leo (gameified caveat (C1)).
\end{activity}

\begin{activity}[title={Activity 8.3 --- Two-analyst roleplay}]
Two students play \emph{Compressor-Anna} (has memorised 4{,}000 deal
memos) and \emph{Understander-Ben} (has read only ten but understands
them). A third student plays the interviewer and asks three
questions about a deal \emph{neither} analyst has seen:

\begin{enumerate}[leftmargin=1.4em,itemsep=2pt]
  \item ``Walk me through how you would price this.''
  \item ``What single fact, if it changed, would flip your pricing
        by more than ten percent?''
  \item ``What is your second-best comparable and why?''
\end{enumerate}

Score each answer for \emph{generalisation} (did the analyst
transfer a principle?) versus \emph{pattern-matching} (did the
analyst return the nearest memo?). Debrief: which analyst is the
LLM in this lecture's framing? (Answer: neither cleanly --- which
is the point of (C5).)

\textbf{Deliverable.} A 3$\times$2 score matrix with notes. Helps
Yuki (concept $\to$ concrete) and Aisha (memorable characters).
\end{activity}

\begin{activity}[title={Activity 8.4 --- Claim-to-caveat matching}]
Match each vendor claim to the caveat that most directly bites it.
More than one match is possible; argue for your primary choice.

\begin{enumerate}[leftmargin=1.4em,itemsep=2pt]
  \item ``Our model is AGI-adjacent.''
  \item ``Our LLM fully understands your firm's legal documents.''
  \item ``Our system approximates the universal Solomonoff prior.''
  \item ``We achieved state-of-the-art perplexity on
        LegalBench-2024.''
  \item ``Our model's reasoning is transparent and auditable.''
\end{enumerate}

Caveats to choose from: (C1) uncomputability, (C2) additive constant
/ small-data regime, (C3) Shannon- vs. Solomonoff-universality,
(C4) necessary vs. sufficient, (C5) extensional vs. intensional.

\textbf{Deliverable.} A table of claim $\to$ primary caveat $\to$
one-sentence rebuttal. Helps Leo (forces re-reading of the formal
paragraphs with a task attached) and Maria (turns philosophy into
operational critique).
\end{activity}

% =========================================================================
